% ===========================================================================
% TEMPLATE.tex — section skeleton for a report in the README.tex format.
% Each section: name + what belongs inside. Replace <fill in> text; keep the
% structure, labels, and \noindent\rule/\FloatBarrier separators. Compile:
%   pdflatex TEMPLATE.tex
% ===========================================================================

\documentclass[12pt,a4paper]{article}
\usepackage[left=2.00cm,right=2.00cm]{geometry}
\usepackage[T1]{fontenc}
\usepackage[utf8]{inputenc}
\usepackage{float}
\usepackage{caption}
\usepackage{graphicx}
\usepackage{longtable}
\usepackage{subcaption}
\usepackage{amsmath}
\usepackage{hyperref}
\hypersetup{colorlinks=true, linkcolor=blue, urlcolor=blue,
            pdftitle={<report title>},
            pdfauthor={<author>}}
\usepackage{placeins}

% ---------------------------------------------------------------------------
% Additions to the supplied preamble:
\usepackage{amssymb}   % \checkmark, which the text uses
\usepackage{booktabs}  % \toprule/\midrule/\bottomrule, consistent table rules
\usepackage{array}     % >{\raggedright\arraybackslash}p{} column specifications
\usepackage[table]{xcolor} % \cellcolor, gradient shading in the cross-case tables
\setlength{\emergencystretch}{2em}   % keeps dense paragraphs within the measure
% ---------------------------------------------------------------------------

\title{<report title>}
\author{<author>}
\date{\today}

\begin{document}
\maketitle

% ===========================================================================
\section*{Introduction}
% ===========================================================================
% One paragraph: what the project set out to do (aim), in one or two sentences.
% Then a short paragraph stating the scale of the work: what was implemented
% (how many methods/variants, over what cases) and how they divide into
% families/groups. End with the headline result (the one-sentence takeaway the
% reader should remember).
%
% <fill in: aim --- e.g. "The aim of this project was to test out different
%  numerical methods for ...">
%
% <fill in: what was implemented and how it divides into families/groups>
%
% <fill in: the headline result --- what wins, what loses, by how much>

% Optional if the text uses abbreviations: a two-column abbreviation table.
\begin{table}[H]
  \centering
  \begin{tabular}{@{}>{\raggedright\arraybackslash}p{0.42\textwidth}
                    >{\raggedright\arraybackslash}p{0.42\textwidth}@{}}
    \toprule
    Full name & Abbreviation \\
    \midrule
    <full name> & <ABBREVIATION> \\
    \bottomrule
  \end{tabular}
\end{table}

\noindent\rule{\textwidth}{0.4pt}

% ===========================================================================
\section*{Table of contents}
% ===========================================================================
% Manual TOC: one \hyperref per numbered section below, in order. Keep it in
% sync with the section list; this report does not use an automatic TOC.
\begin{itemize}
  \item \hyperref[sec:the-experiment]{The experiment}
  \item \hyperref[sec:the-physics]{The physics}
  \item \hyperref[sec:configuration]{Parameters}
  \item \hyperref[sec:the-four-method-families]{The four method families}
  \item \hyperref[sec:analysis]{Analysis --- all fifteen operator combinations}
  \item \hyperref[sec:summary]{Summary}
\end{itemize}

\noindent\rule{\textwidth}{0.4pt}
\FloatBarrier

% ===========================================================================
\section{The experiment}
\label{sec:the-experiment}
% ===========================================================================
% WHAT IS SIMULATED. The physical setup in prose: initial state, the objects
% acting on it, and the variation space (what is chosen independently, how many
% combinations that gives). State the actual parameter values here in passing.
% Structure: setup sentence -> enumerate the terms acting on the system ->
% classify the variants into two structurally different kinds -> the combined
% behaviour in one sentence -> propagation horizon + the extras some methods use
% (as an itemize).
%
% <fill in: initial state and its Bloch vector / purity>
% <fill in: enumerate the acting terms, each with its equation and value>
% <fill in: the two structural kinds of variants and what each does>
% <fill in: combined behaviour, one sentence>
% <fill in: time grid; then itemize the optional add-ons (rates, efficiencies)>

\noindent\rule{\textwidth}{0.4pt}
\FloatBarrier

% ===========================================================================
\section{The physics}
\label{sec:the-physics}
% ===========================================================================
% THE THEORY the numerics implement, in increasing order of specificity:
% general formalism -> the governing equation -> one fully worked example ->
% the conceptual distinctions the variants rest on.

% ---------------------------------------------------------------------------
\subsection{State and Bloch representation}
\label{sec:state-and-bloch-representation}
% ---------------------------------------------------------------------------
% The state space and the representation every method reports through: the
% decomposition equation, how the coefficients are recovered, and purity as a
% function of the vector norm (pure on the sphere, mixed inside).
%
% <fill in: state decomposition equation, recovery rule, purity statement>

% ---------------------------------------------------------------------------
\subsection{The governing equation}
\label{sec:the-governing-equation}
% ---------------------------------------------------------------------------
% The central equation with its terms named and interpreted (one bullet per
% term), the conservation it guarantees, and how independent contributions
% combine (the multi-channel form). Cite the original source.
%
% <fill in: the main equation, labelled>
% <fill in: bullet per term --- what each represents>
% <fill in: trace/conservation statement>
% <fill in: multi-term extension equation>

% ---------------------------------------------------------------------------
\subsection{\texorpdfstring{A worked example: the <name> case}{A worked example: the <name> case}}
\label{sec:this-model-solved-by-hand}
% ---------------------------------------------------------------------------
% One case solved in closed form so the numerics have an exact reference, and a
% note that all other cases follow by permuting operators. Derive the reduced
% equations, state their eigenvalues/timescales, compare the measured numbers
% against the closed form (mark agreement with \checkmark), and give the steady
% state plus the effect of the optional add-on.
%
% <fill in: derivation of the reduced equations>
% <fill in: consequences --- eigenvalues, decay factors, measured-vs-expected>
% <fill in: steady state; effect of the add-on at full strength>

% ---------------------------------------------------------------------------
\subsection{\texorpdfstring{Conceptual axis 1 --- <distinction A>}{Conceptual axis 1 -- <distinction A>}}
\label{sec:conceptual-axis-1}
% ---------------------------------------------------------------------------
% The first independent distinction the variants rest on. Define each side in a
% bullet, state the equivalence the ensemble average forces (the sentence a
% sceptical reader needs), how each side can be implemented, and the special
% value that collapses the distinction.
%
% <fill in: bullet defining each side>
% <fill in: the equivalence statement + both implementation routes>
% <fill in: collapsing value -> what it reduces to>

% ---------------------------------------------------------------------------
\subsection{\texorpdfstring{Conceptual axis 2 --- <distinction B> vs.\ <distinction C>}{Conceptual axis 2 -- <distinction B> vs. <distinction C>}}
\label{sec:conceptual-axis-2}
% ---------------------------------------------------------------------------
% The second independent distinction, as a What changes / single-trajectory /
% ensemble-average table, followed by the measured numbers demonstrating each
% side (itemize).
%
% <fill in: the three-row comparison table>
% <fill in: itemize of measured consequences>

\noindent\rule{\textwidth}{0.4pt}
\FloatBarrier

% ===========================================================================
\section{Parameters}
\label{sec:configuration}
% ===========================================================================
% THE KNOBS. One sentence saying a single set is shared by everything (so all
% comparisons are apples-to-apples), then the Symbol/Value/Meaning table, then a
% derived-quantities paragraph (total time, initial vector, what the run covers
% in units of the natural periods).
%
% <fill in: table --- symbol, value, meaning per parameter>
% <fill in: derived quantities; how many natural periods the run covers>

% ===========================================================================
\section{The four method families}
\label{sec:the-four-method-families}
% ===========================================================================
% THE ALGORITHMS. One \subsection per family (one-paragraph characterisation:
% what it evolves, memory cost, distinctive structural feature), then one
% \subsubsection per method inside it. Every method gets the same recipe:
%   1. the objects it evolves and any precomputed operators (as equations),
%   2. the per-step update (aligned cases/equations),
%   3. what makes it different from its siblings,
%   4. its convergence target and error order,
%   5. its limits --- which special parameter values collapse it onto which
%      other method.

% ---------------------------------------------------------------------------
\subsection{Family 1 --- <name> (role, e.g. deterministic ground truth)}
\label{sec:family-1}
% ---------------------------------------------------------------------------
% One paragraph: how many methods, why they are reproducible, what role they
% play (reference for everyone else).
%
% <fill in: family role paragraph>

\subsubsection{\texorpdfstring{1. <Method name>}{1. <Method name>}}
\label{sec:method-1}
% <fill in: idea, equations, update rule, convergence target, limits>

\subsubsection{\texorpdfstring{2. <Method name>}{2. <Method name>}}
\label{sec:method-2}
% <fill in: idea, equations, update rule, convergence target, limits>

\subsubsection{\texorpdfstring{3. <Method name>}{3. <Method name>}}
\label{sec:method-3}
% <fill in: idea, equations, update rule, convergence target, limits>

\noindent\rule{\textwidth}{0.4pt}
\FloatBarrier

% ---------------------------------------------------------------------------
\subsection{Family 2 --- <name>}
\label{sec:family-2}
% ---------------------------------------------------------------------------
% One paragraph: what the family unravels/evolves, memory cost, structural
% invariant it maintains.
%
% <fill in: family role paragraph>

\subsubsection{\texorpdfstring{4. <Method name>}{4. <Method name>}}
\label{sec:method-4}
% <fill in>

\subsubsection{\texorpdfstring{5. <Method name> $+$ <variant>}{4. <Method name> + <variant>}}
\label{sec:method-5}
% <fill in>

\noindent\rule{\textwidth}{0.4pt}
\FloatBarrier

% ---------------------------------------------------------------------------
\subsection{Family 3 --- <name>}
\label{sec:family-3}
% ---------------------------------------------------------------------------
% One paragraph: what the family evolves, what it permits that the previous
% family cannot, and what it never forms.
%
% <fill in: family role paragraph>

\subsubsection{\texorpdfstring{6. <Method name>}{6. <Method name>}}
\label{sec:method-6}
% <fill in>

\subsubsection{\texorpdfstring{7. <Method name> $+$ <variant>}{7. <Method name> + <variant>}}
\label{sec:method-7}
% <fill in>

\noindent\rule{\textwidth}{0.4pt}
\FloatBarrier

% ---------------------------------------------------------------------------
\subsection{Family 4 --- <name>}
\label{sec:family-4}
% ---------------------------------------------------------------------------
% One paragraph: the operator-sum representation the family is built on, and how
% its discretisation differs structurally from Family 3.
%
% <fill in: family role paragraph>

\subsubsection{\texorpdfstring{8. <Method name>}{8. <Method name>}}
\label{sec:method-8}
% <fill in>

\subsubsection{\texorpdfstring{9. <Method name> $+$ <variant>}{9. <Method name> + <variant>}}
\label{sec:method-9}
% <fill in>

\subsubsection{\texorpdfstring{10. <Method name> $+$ <variant>$\times$2}{10. <Method name> + variant x2}}
\label{sec:method-10}
% <fill in>

\noindent\rule{\textwidth}{0.4pt}
\FloatBarrier

% ---------------------------------------------------------------------------
\subsection{Method comparison at a glance}
\label{sec:method-comparison}
% ---------------------------------------------------------------------------
% A single longtable over all methods: number, name, what it evolves,
% stochastic?, no-jump branch (with error order), convergence target --- and a
% closing sentence naming the two independent routes to the deterministic answer.
%
% <fill in: longtable, one row per method>

% ===========================================================================
\section{Analysis --- <scope of the analysis>}
\label{sec:analysis}
% ===========================================================================

% ---------------------------------------------------------------------------
\subsection{Metrics}
\label{sec:metrics}
% ---------------------------------------------------------------------------
% The closeness metrics, each as a numbered equation (max deviation, MAE, RMSE,
% and a normalised RMSE with its normalisation), how the integrals are
% evaluated, the dimensionless interpretation, and the degenerate-case
% convention. Close by stating what the preceding single-model sections fixed
% and what the sweep varies.
%
% <fill in: metric equations + trapezoid note + null-case convention>
% <fill in: the sweep statement --- what varies, what is held fixed>

% ---------------------------------------------------------------------------
\subsection{Cross-case method agreement}
\label{sec:cross-case-method-agreement}
% ---------------------------------------------------------------------------
% Three full-width metric tables (one per metric) over all cases x all methods,
% cells colour-scaled on one whole-table scale, plus the method-to-reference
% mapping table. Text states: which method is scored against which reference,
% and calls out the degenerate case that is identically zero.
%
% <fill in: reference-mapping table, then the three metric tables>

% ---------------------------------------------------------------------------
\subsection{Case-by-case results}
\label{sec:case-by-case-results}
% ---------------------------------------------------------------------------
% One \subsubsection per case: the operators, the explicit master equation
% integrated, and one figure with the two subfigures (without/with the add-on),
% then \FloatBarrier. Preceded by the observations common to all cases: the
% ranking, the best variant and why, the worst family, and what the control
% column demonstrates.
%
% <fill in: common-observations paragraphs>
%
% --- repeat per case, e.g.: ---
% \subsubsection{\texorpdfstring{<Case heading>}{<Case heading ascii>}}
% \label{sec:case-<slug>}
% \[ <operators> \]
% \[ <explicit equation integrated> \]
% \begin{figure}[H]
%   \centering
%   \begin{subfigure}[b]{0.48\textwidth}
%     \centering
%     \includegraphics[width=\textwidth]{figures/<case>_no_<addon>.png}
%     \caption{Without <addon>}
%   \end{subfigure}\hfill
%   \begin{subfigure}[b]{0.48\textwidth}
%     \centering
%     \includegraphics[width=\textwidth]{figures/<case>_with_<addon>.png}
%     \caption{With <addon>}
%   \end{subfigure}
%   \caption{<Case heading>: <short operator summary>.}
%   \label{fig:<case>}
% \end{figure}
% \FloatBarrier

\noindent\rule{\textwidth}{0.4pt}
\FloatBarrier

% ===========================================================================
\section{Summary}
\label{sec:summary}
% ===========================================================================
% Restate the aim, the scale (methods, families, cases), and the headline
% ranking --- a condensed echo of the Introduction; no new material.
%
% <fill in: aim, scale, headline result>

\end{document}
